\chapter{Apr.~23 --- Complex Analytic Geometry}

\section{Holomorphic Functions}

\begin{remark}
  Let $k = \C$. We will replace polynomials
  with holomorphic functions and consider
  the usual analytic topology on $\C^n$.
\end{remark}

\begin{definition}
  Let $U \subseteq \C$ be open.
  A function $f : U \to \C$ is
  \emph{holomorphic at $c \in U$} (we
  say that $f$ is \emph{holomorphic} if it is
  holomorphic at every point) if one of the
  following equivalent conditions hold:
  \begin{enumerate}
    \item $f$ is \emph{analytic} at $c$;
    \item $f$ is \emph{complex differentiable} at $c$;
    \item the \emph{Cauchy-Riemann
      equations} hold at $c$.
  \end{enumerate}
\end{definition}

\begin{definition}
  Let $U \subseteq \C^n$ be an open set.
  We say that a function
  $f : U \to \C$ is \emph{holomorphic} if
  one of the following equivalent conditions hold:
  \begin{enumerate}
    \item $f$ is analytic, i.e. if at each
      $c = (c_1, \dots, c_n) \in U$,
      there exist $a_J \in \C$ such that
      \[
        f(z)
        = \sum_{J = (j_1, \dots, j_n) \in \N^n}
        a_J (z_1 - c_1)^{j_1} \cdots (z_n - c_n)^{j_n}
      \]
      in some neighborhood around $c$;
    \item at each $c = (c_1, \dots, c_n) \in U$,
      $f(c_1, \dots, c_{i - 1}, z_i, c_{i + 1}, \dots, c_n)$
      is holomorphic in $z_i$ near $c_i$,
    \item the Cauchy-Riemann equations hold:
      Write $f(z) = f(x, y) = u(x, y) + iv(x, y)$, then
      \[
        \frac{\partial u}{\partial x_i}
        - \frac{\partial v}{\partial y_i} = 0
        \quad\text{and}\quad
        \frac{\partial u}{\partial y_i}
        + \frac{\partial v}{\partial x_i} = 0.
      \]
  \end{enumerate}
\end{definition}

\begin{definition}
  A \emph{germ} of a holomorphic function
  at $c \in \C^n$ is a holomorphic
  function $\varphi : U \to \C$ (where
  $c \in U \subseteq \C^n$)
  up to the equivalence relation
  where $(\varphi, U) \sim (\varphi', U')$
  if there exists
  $c \in W \subseteq U \cap U'$ open such
  that $\varphi|_W = \varphi'|_W$. Define
  the space of germs
  \[
    \OO_n
    = \{
      \text{germs of holomorphic functions at $0 \in \C^n$}
    \}.
  \]
\end{definition}

\begin{prop}\label{prop:properties-of-On}
  We have the following:
  \begin{enumerate}
    \item $\OO_n$ is a ring;
    \item there is an inclusion
      $\OO_n \to \C \llbracket z_1, \dots, z_n \rrbracket$
      (we often write $\C \langle z_1, \dots, z_n \rangle$ for the image);
    \item $\OO_n$ has a unique
      maximal ideal
      $\m = \{\varphi \in \OO_n : \varphi(0) = 0\}$
      (so $\OO_n$ is a local ring);
    \item $\OO_n$ is Noetherian;
    \item $\OO_n$ is a UFD.
  \end{enumerate}
\end{prop}

\begin{remark}
  When $n = 1$, the ideals of $\OO_1$
  are of the form $(z^m)$ with $m \ge 0$
  and $(0)$. In general, proving
  Proposition \ref{prop:properties-of-On}(4)-(5)
  is more difficult and relies on the
  \emph{Weierstrass division theorem}.
\end{remark}

\begin{remark}
  If we take an open set $U \subseteq \C^n$,
  we can consider
  \[
    \OO_{\C^n}(U)
    = \{\text{holomorphic functions $U \to \C$}\}.
  \]
  However, $\OO_{\C^n}(U)$ is not
  necessarily Noetherian. For example,
  in $\OO_\C(\C)$ we can take
  \[
    I_k = \{
      \varphi \in \OO_\C(\C)
      : f(i) = 0 \text{ for all } i \in \N
      \text{ with } i \ge k
    \}.
  \]
  Then $I_k \subseteq I_{k + 1} \subseteq I_{k + 2} \subseteq \cdots$
  and the chain is strict: Something like
  $\sin(\pi z) / (z - k)$ lies in
  $I_{k + 1} \setminus I_k$.
\end{remark}

\section{Analytic Subvarieties}
\begin{definition}
  Let $U \subseteq \C^n$ be an open
  set. An \emph{analytic subvariety}
  of $U$ is a subset $X \subseteq U$ such
  that for each $x \in X$, there exists an
  open neighborhood $x \in V \subseteq U$
  such that
  \[
    V \cap X
    = \{x \in V : f_1(x) = \cdots = f_r(x) = 0\}
    =: V(f_1, \dots, f_r)
  \]
  for some $f_1, \dots, f_r \in \OO_{\C^n}(V)$.
\end{definition}

\begin{example}
  The following are analytic subvarieties:
  \begin{enumerate}
    \item Hypersurfaces:
      Given $f \in \OO_{\C^n}(U)$, we have
      $V(f) = \{x \in U : f(x) = 0\}$.
    \item $Z = V(w - e^z) \subseteq \C^2$.
      Note that $Z \cap V(w - 1) = \{(1, 2\pi i k) : k \in \Z\}$,
      which is notably infinite.\footnote{We saw that with polynomials, the intersection of two distinct curves is finite.}
    \item Let $U = \C_{w, z}^2 \setminus \{w = 0\}$
      and $X = \bigcup_{n = 1}^\infty V(z - nw)$.
      Note that
      $\overline{X} \subseteq \C^2$
      is \emph{not} an analytic subvariety of $\C^2$.
      This is intuitively because
      $\codim_{\C^2} \{w = 0\} = 1$ and
      $\codim_U X = 1$.
  \end{enumerate}
\end{example}

\begin{definition}
  A \emph{germ} of a complex analytic
  subvariety of $\C^n$ at $0$ is a complex
  analytic subvariety
  $0 \in X \subseteq U \subseteq \C^n$
  (where $U$ is open in $\C^n$), up to the
  equivalence relation where
  $X \sim X'$ if there exists
  $0 \in W \subseteq U \cap U'$ such that
  $X \cap W = X' \cap W$.
\end{definition}

\begin{remark}
  Given a germ $0 \in X \subseteq \C^n$ of
  a complex analytic subvariety, we can
  define
  \[
    I(X)
    = \{
      \varphi \in \OO_n :
      \text{$\varphi$ vanishes on a neighborhoods of $X$ containing $0$}
    \}.
  \]
  Similarly, given an ideal $I \subseteq \OO_n$,
  we can define
  \[
    0 \in X = V(f_1, \dots, f_r),
  \]
  where $f_1, \dots, f_r$ are generators
  of $I$ (recall that $\OO_n$ is Noetherian)
  defined on some open neighborhood of
  $0 \in U \subseteq \C^n$.
  Note that $X$ itself is not an
  analytic subvariety, but a germ
  of an analytic subvariety.
\end{remark}

\begin{prop}
  We have the following:
  \begin{enumerate}
    \item if $X_1 \subseteq X_2$ are
      germs of analytic subvarieties
      (i.e. there is some neighborhood
      $0 \in U \subseteq \C^n$ where
      $X_1 \cap U \subseteq X_2 \cap U$), then
      $I(X_1) \supseteq I(X_2)$;
      \pagebreak
    \item if $I_1 \subseteq I_2 \subseteq \OO_n$
      are ideals, then
      $V(I_1) \supseteq V(I_2)$;
    \item $V(I(X)) = X$;
    \item $I(V(I)) = \sqrt{I}$;
    \item $X$ is irreducible if and only if
      $I(X)$ is prime;
    \item any germ of a complex analytic
      subvariety at $0 \in \C^n$
      can be written uniquely (up to
      reordering) as a union of
      irreducible ones.
  \end{enumerate}
\end{prop}

\begin{remark}
  Consider the curve
  $x^2 - y^2 - y^3 = 0$. Recall that
  this is irreducible
  in algebraic geometry (i.e. in the
  Zariski topology). However, in complex
  analytic geometry, we can write
  \[
    x^2 - y^2 - y^3
    = x^2 - y^2(1 + y)
    = (x - y\sqrt{1 + y})(x + y\sqrt{1 + y})
  \]
  for $|y| < 1$, so this curve is
  not irreducible as a germ of a
  complex analytic subvariety.
\end{remark}

\begin{remark}
  The finiteness of the irreducible
  decomposition fails for non-germs:
  Consider
  \[
    \{0, 1, 2, \dots\} \subseteq \C.
  \]
\end{remark}

\section{Complex Manifolds}

\begin{definition}
  A \emph{complex manifold} is a
  locally ringed space $(X, \OO_X)$ such that
  \begin{itemize}
    \item $X$ is Hausdorff with countable
      basis;
    \item $\OO_X$ is a subsheaf of the
      sheaf of $\C$-valued continuous
      functions on $X$;
    \item at any $x \in X$, there exists
      $x \in U \subseteq X$ open such that
      \[
        (U, \OO_X|_U) \cong (V, \OO_{\C^n}|_V)
      \]
      for some open set $0 \in V \subseteq \C^n$.
  \end{itemize}
\end{definition}

\begin{example}
  The following are examples of complex
  manifolds:
  \begin{enumerate}
    \item $X = \C^n$ or $\C \PP^n$.
    \item $X = V(f) \subseteq \C^n$
      such that $f \in \OO_{\C^n}(\C^n)$
      and $[\partial f / \partial z_1, \cdots, \partial f / \partial z_n] \ne 0$
      on $X$.
  \end{enumerate}
\end{example}

\begin{definition}
  A \emph{complex analytic closed subvariety}
  of a complex manifold $X$ is a closed
  set $Z \subseteq X$ that is locally a
  complex analytic subvariety for an open
  ball $V \subseteq \C^n$.
\end{definition}

\begin{theorem}[Chow's theorem]
  If $X \subseteq \C \PP^n$ is an analytic
  subvariety, then there exist
  homogeneous polynomials
  $F_1, \dots, F_r \in \C[z_0, \dots, z_n]$ 
  such that $X = V(F_1, \dots, F_r)$.
\end{theorem}

\begin{proof}
  Consider the quotient map
  $q : \C^{n + 1} \setminus \{0\} \to \C \PP^n$.
  Note that $Z = q^{-1}(X)$ is a complex
  analytic subvariety of $\C^{n + 1} \setminus \{0\}$.
  Now we have
  \[
    \codim_{\C^{n + 1}} \{0\}
    = n + 1 > \codim_{\C^{n + 1} \setminus \{0\}} Z,
  \]
  so Levi's extension theorem says that
  $\overline{Z}$ is a complex analytic
  subvariety of $\C^{n + 1}$.
  Let $I = I(\overline{Z}) \subseteq \OO_{n + 1}$.
  Now for any $f \in I$, we can write
  $f(z) = \sum_{j \ge 0} f_j(z)$ with
  $f_j$ a homogeneous polynomial
  of degree $j$.

  For any fixed $a \in \overline{Z}$, we have
  \[
    0 = f(\lambda a)
    = \sum_{j \ge 0} \lambda^j f_j(a)
  \]
  for every $\lambda \in \C$, so
  $f_j(a) = 0$ for all $j \ge 0$. So
  $f_j \in I$.  Thus there are
  homogeneous polynomials $(F_j)_{j \in J}$
  that generate $I$ (up to infinite sums).
  This gives that
  $V(F_j : j \in J) = X$, and
  we can then choose a finite collection
  cutting out $X$ by Noetherianity.
\end{proof}

\begin{example}
  Consider $X = \{1, 2, \dots, \} \subseteq \C^1 \subseteq \C \PP^1$.
  Now $X$ is an analytic subvariety of
  $\C^1$, but it is not an analytic
  subvariety of $\C \PP^1$.
\end{example}
